%% P10 \dash{} Rozkład SVD, aproksymacja niskiego rzędu i uwarunkowanie
%%
%% Zalążek. Poniższy opis jest kontraktem tego programu: co ma uzasadnić i co
%% ma dać czytelnikowi. Napisanie programu oznacza usunięcie bloku
%% \programstub{} i zastąpienie go programem. Jeśli po przerobieniu materiału
%% opis okaże się błędny, popraw go i odnotuj sprzeczność w CLAUDE.md w sekcji
%% *Resolved questions*.

\program{Rozkład SVD, aproksymacja niskiego rzędu i uwarunkowanie}
\label{prog:P10}

\programstub{%
\textit{[Opis do przetłumaczenia \dash{} patrz wydanie angielskie.]}

Argument: every matrix, square or not, factors into rotate--stretch--rotate.
Truncating the stretch gives the provably best approximation of a given rank
(Eckart--Young, stated). The ratio of largest to smallest singular value is
the condition number, and it is the amplification factor from input error to
output error. Payoff: the singular-value spectrum of a real embedding
matrix, plotted, showing how few directions carry the energy \dash{} the
empirical case for LoRA and for embedding compression, measured; the
pseudoinverse as least squares done properly; why the normal equations
square the condition number and a QR solve does not, which is the concrete
form of \enquote{do not invert}; and the condition number as the number that
predicts how many iterations an optimiser will need, handed forward to P19.

\medskip
\textbf{Planowana długość:} 60 ramek.
}
